{\small
\textbf{Abstract.} A retail world model is meant to answer ``what if'': what sales a plan the merchants have never
run would bring. We test whether it does, in generated stores where the true effect of every plan is known, and
what fixes it when it does not. Models are trained on the plans merchants make and asked to forecast plans held out
by design: deeper cuts than any in training, lasting regular-price moves, promotions moved to quiet weeks and no
plan at all. Across \nFits{} fitted models we find four things. First, forecast accuracy is a fair guide between
kinds of model (rank correlation with error on new plans \linkAllMin{} to \linkAllMax{}) and a poor guide between
versions of one world model: where merchants plan into high-demand weeks the link reverses, down to \linkBWmOne{}
(interval \linkBWmOneCi{}). Second, training the world model to learn the effect of actions (JEPA pretraining, a
latent consistency check, adaptation at the plan date) leaves its error on new plans unchanged (\errLActLo{} to
\errLActHi{} against \errLWm{}). Third, what fixes it is the recorded regular price and three price rules: sales may
not fall as a cut deepens, the gain from deeper cuts may flatten, and a lasting regular price acts through its own
elasticity. Given the same rules, a log-linear model is best when its curve can match the truth (error \errBLlShape{}
against \errBShape{} for the world model with the rules and \errBLgbShape{} for LightGBM with the rules); the world
model with the rules is best when the true curve lies outside them (\errKShape{} against \errKLlShape{} and
\errKLgbShape{}, ahead of the log-linear model in \shapeBeatsLlK{} of twelve worlds). Fourth, a small share of random
test plans helps every model a little. For retailers: give the model the recorded regular price and sensible price
rules, use a simple model when the shape of the price response is known and a world model when it is not, and
choose versions on held-out plans rather than on forecast accuracy.
\par}
